\section{Failed Approaches: Experimental Journey to the Final Design}
\label{sec:failed-approaches}

The final architecture and loss described in this chapter emerged from a
sequence of experiments, most of which failed.
This section documents those dead ends in the order they were encountered,
explaining \emph{why} each approach failed and what each failure revealed.
Several of the failures also point directly to constraints on future extensions.

% ============================================================
\subsection{Phase 1: Displacement Regression with L1 Loss}

\paragraph{Approach.}
The first attempt was the most natural formulation: train a sparse U-Net to
predict per-point displacement vectors $\Delta_i$ directly, supervised by L1
loss against precomputed nearest-neighbour targets
$d_{gt}(\hat{p}_i) = p^*_{nn}(\hat{p}_i) - \hat{p}_i$.

\paragraph{Failure: Zero-inflation collapse.}
As analysed in Section~\ref{sec:loss}, 60\% of decoded points at r2 have
$d_{gt} = \mathbf{0}$ (they are already correctly placed).
The gradient of the L1 loss is $\pm 1$ scaled by the target magnitude.
When 60\% of gradients push toward zero and 40\% push outward, the mean
gradient across a mini-batch is biased strongly toward zero.
Within 5--10 epochs, every network trained with this loss collapses to
predicting near-zero displacements for all points, regardless of architecture
size, learning rate, or regularisation.
By epoch 10, the refined cloud is effectively unchanged from the input,
producing PSNR improvements of $< 0.05$~dB.

\paragraph{Attempted mitigations.}
\begin{itemize}
  \item \textit{Smooth-L1 (Huber) loss}: replaces L1 with L2 near zero and
    L1 far from zero. This changes the gradient scale but not the directional
    bias. Zero-collapse still occurred, only slightly slower.

  \item \textit{Curvature-weighted L1}: weight each term by $1 + \alpha\kappa_i$,
    amplifying gradient from high-curvature points (edges).
    This helps the non-zero gradient terms dominate at edges,
    but flat-region zero-inflation still overwhelms the total gradient.
    Collapse to near-zero for flat points, near-correct for edge points:
    a partial improvement but not a solution.

  \item \textit{Min-of-$K$ target selection}: instead of using the 1-NN target,
    evaluate $K = 5$ candidate targets (nearest 5 neighbours in $\pc^*$) and
    take the one with minimum squared distance to the predicted displacement.
    This addresses Problem~2 (NN noise) but leaves Problem~1 (zero-inflation)
    entirely intact.
\end{itemize}

None of these mitigations resolved the problem.
The root cause is structural: a regression loss against a target distribution
with a heavy mass at zero is biased toward the zero predictor, regardless of
architecture or weighting.

% ============================================================
\subsection{Phase 2: Gated Architecture (Move / No-Move Prediction)}

\paragraph{Approach.}
To directly address zero-inflation, we added a binary gate output to the network:
a per-point probability $g_i \in [0, 1]$ predicting whether the point should
be moved ($g_i > 0.5$) or kept in place ($g_i < 0.5$).
Two gated architectures were implemented, inspired by learned gating
mechanisms for sparse prediction~\cite{yu2019freeform}:
\begin{itemize}
  \item \textit{GatedSparseUNet}: shared backbone, separate gate head and
    displacement head, joint training with binary cross-entropy + L1 loss.
  \item \textit{ThresholdGatedUNet}: train gate first, then freeze it and
    train displacement head on predicted-to-move points only.
\end{itemize}

\paragraph{Failure: Gate does not decouple the problem.}
The gate head learns to predict which points have non-zero targets, which
requires learning the same geometric information as the displacement head.
This means the gate effectively learns a noisy version of the displacement
magnitude, and its cross-entropy loss has the same zero-inflation issue:
if 60\% of points are labeled ``no-move'' (class 0), the gate collapses to
predicting class 0 everywhere.
The two-stage ThresholdGated variant improves on this but suffers from poor
gate quality during stage 1, causing stage 2 to train on a biased subset.

\paragraph{Marginal improvement.}
The gated models achieved small PSNR gains ($+0.1$--$0.2$~dB at r2), suggesting
that the gating structure is directionally correct but insufficient.
The gate is solving the wrong problem: trying to classify points
rather than treating the zero-inflation as a \emph{distribution problem}
that requires a different loss.

% ============================================================
\subsection{Phase 3: Stratified Loss (Laplacian on Edges, Zero-Penalty on Flat)}

\paragraph{Approach.}
Rather than a single loss for all points, apply different losses to different
geometric regions:
\begin{itemize}
  \item For high-curvature points: a Laplacian smoothness loss penalising
    non-smooth displacements.
  \item For low-curvature points: an explicit L2 penalty on the displacement
    magnitude (push displacements to zero for flat points).
\end{itemize}

The motivation was to directly encode the two failure modes into the loss.

\paragraph{Failure: Two problems treated as one.}
This approach addresses Problems~1 and~2 together rather than independently.
The Laplacian loss on edge points adds a regularity constraint but provides no
target supervision: it penalises rough displacements without specifying
\emph{where} to move.
The zero-penalty on flat points is directionally correct but leaves edge
displacement without adequate supervision.

Curvature estimated on the decoded cloud is also unreliable at edge regions
(the decoded cloud is smooth there by construction).
The stratification boundaries are therefore noisiest precisely where they matter
most.

The stratified loss produced moderate improvement ($+0.3$--$0.5$~dB at r2)
but was sensitive to the curvature threshold hyperparameter and did not
generalise well across sequences.

% ============================================================
\subsection{Phase 4: Laplacian Preservation Loss}

\paragraph{Approach.}
Inspired by geometry processing~\cite{sorkine2004laplacian} and its use
as a deep learning regulariser~\cite{wang2018pixel2mesh}, we tried adding a
Laplacian preservation term: the refined cloud should have the same Laplacian
structure as the original cloud.
Two implementations were tested:
\begin{itemize}
  \item \textit{kNN Laplacian}~\cite{taubin1995signal}: build a $k$-NN graph
    on the decoded cloud, compute Laplacian coordinates, and penalise
    deviations of the refined coordinates from target Laplacian coordinates.
  \item \textit{Voxel Laplacian via sparse convolution}: implement the
    Laplacian as a $3 \times 3 \times 3$ sparse convolution with fixed
    centre-minus-mean kernel, avoiding dynamic graph construction.
\end{itemize}

\paragraph{Failure: Laplacian amplifies noise.}
The Laplacian operator is high-pass: it amplifies differences from local means.
On a decoded cloud with oversmoothing artifacts, the Laplacian of the decoded
cloud is already corrupted.
Penalising deviations from this corrupted target propagates the oversmoothing
bias into the displacement predictions.

Additionally, the kNN Laplacian requires building a dynamic graph at every
training step, which is prohibitively slow ($\sim$10$\times$ training cost
vs.\ standard loss).
The voxel Laplacian is fast but uses a fixed neighbourhood structure (the
$3^3$ grid) that is not adapted to the local point density.

Neither variant improved over the baseline L1 loss in ablation experiments.
The Laplacian approach was abandoned after these negative results.

% ============================================================
\subsection{Phase 5: Combined Chamfer Distance + Laplacian}

\paragraph{Approach.}
After discovering that the Chamfer Distance loss resolves zero-inflation
(see Section~\ref{ssec:dcd}), we attempted to combine it with a Laplacian
smoothness term:
\begin{equation}
  \mathcal{L} = \mathcal{L}_{\text{CD}} + \lambda \mathcal{L}_{\text{Lap}},
\end{equation}
hypothesising that the Laplacian term would prevent the CD loss from
over-correcting (moving points to positions not on the surface).

\paragraph{Failure: Laplacian is unnecessary and harmful.}
The symmetric Chamfer Distance already prevents over-correction via its
reverse term ($\pc^*_{\text{crop}} \to \hat{\pc}_{\text{ref}}$): if the refined
cloud is pulled away from the surface, the reverse term increases, providing
a corrective gradient.
The Laplacian adds a conflicting regularisation that prevents large but correct
displacements at edge regions.
In practice, the combined loss produced smaller PSNR gains at r2 than CD alone
($+1.1$~dB vs.\ $+1.45$~dB), and the tuning of $\lambda$ was unstable.

The conclusion was that the Chamfer Distance is self-regularising and no
auxiliary term is needed.

% ============================================================
\subsection{Phase 6: Deep FiLM Conditioning (Rate on All ResBlocks)}

\paragraph{Approach.}
The first multi-rate extension applied FiLM conditioning to every ResBlock
of both encoder and decoder (architecture: FiLMSparseUNet, v3 and variants).
The rate scalar $r = \log_2(\bpp)$ was fed through a small MLP, and the
resulting $\gamma$ and $\beta$ vectors were broadcast to modulate features
at each of the 6 ResBlocks.

\paragraph{Failure: Gradient conflict in shared backbone.}
When training on r2 ($\bpp = 0.05$), r4 ($\bpp = 0.15$), and r6 ($\bpp = 0.30$)
simultaneously with deep FiLM:
\begin{itemize}
  \item The r2 gradient signal says: ``apply large displacements to edge points.''
  \item The r6 gradient signal says: ``apply near-zero displacements everywhere.''
\end{itemize}
These gradients conflict in the shared backbone parameters.
Because the FiLM $\gamma/\beta$ parameters also appear in the backbone, the
conflict propagates through the conditioning pathway as well.

The symptom was a flat forward Chamfer Distance loss plateau at $\approx 2.45$
for 12 consecutive epochs, indicating that the backbone was unable to make
progress in any direction.
This is a gradient cancellation phenomenon: each batch of patches from r2 and
r6 produces opposing gradient updates that largely cancel.

\paragraph{Fix: Rate-invariant backbone + head-only FiLM.}
The solution was to remove FiLM from the backbone entirely, making it
rate-invariant, and apply FiLM conditioning only to the final 32-channel
pre-head features.
The backbone learns geometry-invariant features that are useful at all rates;
the FiLM head modulates these features to produce rate-appropriate displacement
magnitudes.
This is the \textit{FiLMHeadSparseUNetV2} architecture described in
Section~\ref{sec:film-head}, which resolved the gradient conflict and achieved
the best multi-rate results.

% ============================================================
\subsection{Phase 7: Chamfer Boundary Padding}

\paragraph{Approach.}
During Chamfer Distance implementation, we initially added a boundary padding
of $d = 10$ voxels to the original cloud crop:
\begin{equation}
  \pc^*_{\text{crop}} = \{p^* \in \pc^* : p^* \in
    [u_{\min} - d,\, u_{\max} + d]^3\}.
\end{equation}
The motivation was to give boundary decoded points access to nearby original
points that fall just outside the patch bounding box.

\paragraph{Failure: Reverse term dominates.}
With $d = 10$, the crop $\pc^*_{\text{crop}}$ is significantly larger than
$\hat{\pc}_{\text{patch}}$.
Boundary decoded points near the edge of the patch find their nearest original
point at distance $\approx d$ voxels (outside the original patch boundary).
The reverse term (original points to refined cloud) then contains many points
far from any refined point, dominating the total loss.

By epoch 3, the reverse term constitutes over 99\% of the total loss.
The effective gradient tells the network: ``move all decoded points outward
toward the patch boundary,'' which is a degenerate signal that prevents any
geometric learning.

Setting $d = 0$ (no boundary padding) fixes the problem.
With zero padding the two terms are balanced, and the model learns correct
displacements within 3--5 epochs.

% ============================================================
\subsection{Phase 8: Adaptive Rate Weighting (Kendall and Dynamic)}

\paragraph{Approach.}
Rather than manually setting rate weights $[w_2, w_4, w_6]$, we tried two
principled alternatives from the multi-task learning literature:

\begin{itemize}
  \item \textit{Kendall uncertainty weighting~\cite{kendall2018multi}}: learns
    per-task log-uncertainty parameters $\log\sigma_k$; the effective weight
    is $w_k = 1/(2\sigma_k^2)$.

  \item \textit{Dynamic detached weighting}: at each epoch, set
    $w_k \propto \mathcal{L}_k / \bar{\mathcal{L}}$ (proportional to current
    loss normalised by mean), clamped to $[0.1, 10.0]$.
\end{itemize}

\paragraph{Failure: Mismatch between rate learning dynamics.}
The Kendall method failed catastrophically: within 5 epochs, the weights
collapsed to the inverted ordering $w_2 < w_4 < w_6$.
The reason is that Kendall's homoscedastic uncertainty is designed for tasks
with genuinely different noise levels.
Here, the r2 task has a larger Chamfer Distance loss (more displacement required)
which Kendall interprets as higher uncertainty, causing it to \emph{downweight}
the most important rate.
This is the opposite of the correct inductive bias.

The dynamic detached method converged to a nearly static ratio
$\approx [1.27, 0.91, 0.82]$ by epoch 5.
This ratio reflects the proportional learning progress across rates: since the
shared backbone causes all rates to improve at proportional rates, the dynamic
normaliser quickly stabilises.
The best r2 edge improvement with dynamic weighting was $+1.28$~dB, compared
to $+2.06$~dB with static weights.

Adaptive weighting helps when tasks have decoupled learning dynamics, i.e.,
when separate parameters handle each task.
Here the shared backbone ties the rates together, so the static weights
$[1.0, 0.3, 0.05]$, which encode the prior that r2 has the richest signal,
cannot be recovered by any data-driven scheme applied to the shared loss.

% ============================================================
\subsection{Phase 9: Post-Hoc Curvature Scaling}

\paragraph{Approach.}
After the Chamfer Distance model achieved $+1.45$~dB at r2, we attempted to
further boost edge improvements by post-hoc scaling: multiply the predicted
displacement by $\alpha \kappa_i$ where $\kappa_i$ is the estimated surface
curvature of point $\hat{p}_i$.
The hypothesis was that the model's predictions at edge points would be amplified
appropriately, recovering additional PSNR at the cost of flat-region quality.

\paragraph{Failure: Model already correct.}
For all scaling factors $\alpha > 1.0$ tested (1.2, 1.5, 2.0, 3.0), the D1
PSNR decreased compared to no scaling.
The reason: the model had already learned to apply larger displacements at
edge regions (it observes the local coordinate structure and implicitly infers
curvature).
Post-hoc scaling on top of already-correct predictions produces
\emph{over-correction}: edge points are pushed past the target surface,
increasing the D1 error.

The failure of post-hoc scaling was the first evidence that the network was
already learning curvature-adaptive behaviour from the coordinate features
alone, which motivated the curvature input ablation (Section~\ref{ssec:ablation-curvature}).

% ============================================================
\subsection{Summary of Key Lessons}

Table~\ref{tab:failed-approaches} summarises the failed approaches and the
lesson each contributed to the final design.

\begin{table}[t]
  \centering
  \caption{%
    Summary of failed approaches and their contributions to the final design.
    ``Gain'' is the approximate r2 edge PSNR improvement over the unrefined baseline.
  }
  \label{tab:failed-approaches}
  \small
  \begin{tabular}{llp{5.5cm}}
    \toprule
    \textbf{Approach} & \textbf{r2 Edge Gain} & \textbf{Key Lesson} \\
    \midrule
    L1 displacement regression & $< 0.05$~dB & Zero-inflation collapse is structural \\
    Smooth-L1 & $< 0.05$~dB & Gradient bias unchanged \\
    Curvature-weighted L1 & $\approx 0.1$~dB & Partial but insufficient fix \\
    Min-of-$K$ targets & $\approx 0.1$~dB & Fixes NN noise but not zero-collapse \\
    Gated architecture & $\approx 0.15$~dB & Gate collapses for same reason \\
    Stratified loss & $\approx 0.4$~dB & Conflates two independent problems \\
    Laplacian preservation & $\approx 0.0$~dB & Laplacian of decoded cloud is corrupt \\
    CD + Laplacian & $\approx 1.1$~dB & CD is already self-regularising \\
    Deep FiLM on all ResBlocks & plateau & Gradient conflict in shared backbone \\
    CD with boundary padding=10 & negative & Reverse term dominates; degenerate gradient \\
    Kendall rate weighting & worse & Inverts the correct inductive bias \\
    Post-hoc curvature scaling & negative & Model already rate-curvature adaptive \\
    \midrule
    \textbf{Final: CD + head FiLM + static weights} & $\mathbf{+2.06}$~dB & -- \\
    \bottomrule
  \end{tabular}
\end{table}

The progression from $< 0.05$~dB (L1 collapse) to $+2.06$~dB (final model)
is entirely explained by identifying and resolving three independent failure
modes: (1) zero-inflation in the loss, (2) gradient conflict in the architecture,
and (3) boundary artifacts in the patch sampling.
No single change was sufficient; each required the previous fix to be in place
before its benefit could be observed.

% ============================================================
\begin{table}[t]
  \centering
  \caption{%
    Summary of the final post-processing model (v10-nocurv, epoch 24 checkpoint).
  }
  \label{tab:model-summary}
  \small
  \begin{tabular}{ll}
    \toprule
    \textbf{Component} & \textbf{Configuration} \\
    \midrule
    \multicolumn{2}{l}{\textit{Input}} \\
    Features & 3-channel normalised $(x, y, z)$ coordinates \\
    Patch size & $64^3$ voxels, stride 32 (50\% overlap) \\
    Augmentation & 48 cube symmetries (uniform random) \\
    \midrule
    \multicolumn{2}{l}{\textit{Backbone}} \\
    Architecture & 3-level sparse U-Net (MinkowskiEngine) \\
    Channels & 32 / 64 / 128 (encoder); 128 / 64 / 32 (decoder) \\
    Kernel size & $3^3$ submanifold sparse convolution \\
    Downsampling & Stride-2 sparse convolution (\textit{E}$_2$, \textit{E}$_3$) \\
    Skip connections & Feature concat at resolutions 1$\times$, 2$\times$ \\
    \midrule
    \multicolumn{2}{l}{\textit{Rate Conditioning}} \\
    Rate input & $r = \log_2(\bpp)$, scalar \\
    Rate MLP & Linear(1,64) + ReLU + Linear(64,64) \\
    FiLM output & $\boldsymbol{\gamma} \in \R^{32}$, $\boldsymbol{\beta} \in \R^{32}$ \\
    Modulation & Channel-wise: $\mathbf{h}' = \boldsymbol{\gamma} \odot \mathbf{h} + \boldsymbol{\beta}$ \\
    \midrule
    \multicolumn{2}{l}{\textit{Output}} \\
    Head & SparseConv $1^3$, then $5 \cdot \tanh(\cdot)$ \\
    Output & Per-point displacement $\Delta_i \in [-5, 5]^3$ \\
    \midrule
    \multicolumn{2}{l}{\textit{Loss}} \\
    Loss function & Symmetric Chamfer Distance, chamfer\_padding=0 \\
    Rate weights & $[w_2, w_4, w_6] = [1.0, 0.3, 0.05]$ \\
    \midrule
    \multicolumn{2}{l}{\textit{Training}} \\
    Optimiser & AdamW ($\eta = 10^{-3}$, $\lambda = 10^{-4}$) \\
    Schedule & 3-epoch warm-up + cosine annealing, 30 epochs \\
    Batch size & 8 patches \\
    Total parameters & $\approx 2.5\text{M}$ \\
    Best checkpoint & Epoch 24 (validation loss) \\
    \bottomrule
  \end{tabular}
\end{table}
